\documentclass[12pt]{article}
\usepackage[margin=1in]{geometry}
\usepackage[utf8]{inputenc}
\usepackage{latexsym, amsfonts, enumitem, graphicx, environ,enumitem, fancyhdr, color, amssymb, amsthm, amsmath, mathtools, tikz, nicematrix}
\usetikzlibrary{patterns}
\usetikzlibrary{decorations.pathreplacing}
\pagestyle{fancy}
\setlength{\headheight}{65pt}
\newenvironment{problem}[2][Problem]{\begin{trivlist}
    \item[\hskip \labelsep {\bfseries #1}\hskip \labelsep {\bfseries #2.}]}{\end{trivlist}}
\DeclareMathOperator{\Ima}{Im}

\usepackage{hyperref}
\usepackage[capitalize]{cleveref}

\newtheorem{thm}{Theorem}[section]
\newtheorem{lem}[thm]{Lemma}
\newtheorem{cl}[thm]{Claim}
\newtheorem{prop}[thm]{Proposition}
\newtheorem{cor}[thm]{Corollary}
\newtheorem{conj}[thm]{Conjecture}
\newtheorem{obstruction}[thm]{Obstruction}

\usepackage{mdframed}

\theoremstyle{definition}
\newtheorem{definition}[thm]{Definition}
\newtheorem{remark}[thm]{Remark}
\newtheorem{question}{Question}
\newtheorem{example}[thm]{Example}
\newtheorem{exercise}[thm]{Exercise}
\newtheorem{properties}[thm]{Properties}

\usepackage{tcolorbox}
\tcbuselibrary{theorems,skins,breakable}

\tcolorboxenvironment{example}{
enhanced jigsaw,colframe=cyan,interior hidden,
breakable,%before skip=10pt,after skip=10pt
}

\tcolorboxenvironment{thm}{
enhanced jigsaw,colframe=red,interior hidden,
breakable,%before skip=10pt,after skip=10pt
}

\tcolorboxenvironment{prop}{
enhanced jigsaw,colframe=red,interior hidden,
breakable,%before skip=10pt,after skip=10pt
}

\tcolorboxenvironment{properties}{
enhanced jigsaw,colframe=red,interior hidden,
breakable,%before skip=10pt,after skip=10pt
}

\tcolorboxenvironment{lem}{
enhanced jigsaw,colframe=red,interior hidden,
breakable,%before skip=10pt,after skip=10pt
}
\tcolorboxenvironment{cor}{
enhanced jigsaw,colframe=red,interior hidden,
breakable,%before skip=10pt,after skip=10pt
}
\tcolorboxenvironment{obstruction}{
enhanced jigsaw,colframe=red,interior hidden,
breakable,%before skip=10pt,after skip=10pt
}

\tcolorboxenvironment{proof}{
enhanced jigsaw, frame hidden,%interior hidden,
breakable,%before skip=10pt,after skip=10pt
}

\lhead{Avi Chetlin \\ Assignment 01 \\ 31 August 2026}
\rhead{Dr. David Singer \\ MATH 465 Differential Geometry \\ Fall 2026}

\begin{document}
\begin{problem}{1.2.5}
	Let \(\alpha: I \rightarrow \mathbb{R}^3\) be a parametrized curve, with \(\alpha^{\prime}(t)\neq 0\) for all \( t \in I \). Show that \( \left| \alpha(t) \right| \) is a nonzero constant if and only if \(\alpha(t)\) is orthogonal to \(\alpha^{\prime}(t)\) for all \( t \in I \).
\end{problem}
\begin{proof}
	Let \(\alpha: I \rightarrow \mathbb{R}^3\) be a regular parametrized curve, and write
	\begin{equation}
		\label{eq:3}
		\alpha(t) = (x(t), y(t), z(t)).
	\end{equation}
	By definition,
	\begin{equation}
		\label{eq:2}
		\left| \alpha(t) \right|^2 = x(t)^2 + y(t)^2 + z(t)^2,
	\end{equation}
	hence
	\begin{equation}
		\label{eq:4}
		\frac{\mathrm{d} }{\mathrm{d} t}\left| \alpha(t) \right|^2 = 2(x(t)x^{\prime}(t) + y(t)y^{\prime}(t) + z(t)z^{\prime}(t)).
	\end{equation}

	Now, \( \left| \alpha(t) \right| \) is a nonzero constant if and only if \( \frac{\mathrm{d} }{\mathrm{d} t}\left| \alpha(t) \right|^2 = 0\) and \( \left| \alpha(t) \right|^2 \neq 0\) for all \( t \in I \).
	That is,
	\begin{equation}
		\label{eq:5}
		x(t)^2 + y(t)^2 + z(t)^2 \neq 0,
	\end{equation}
	and by \cref{eq:4},
	\begin{equation}
		\label{eq:6}
		x(t)x^{\prime}(t) + y(t)y^{\prime}(t) + z(t)z^{\prime}(t) = 0
	\end{equation}
	for all \( t \in I \).

	Noting that \cref{eq:6} is, by definition, the dot product \(x(t) \cdot x^{\prime}(t)\), we have that \( \left| \alpha(t) \right| \) is a nonzero constant if and only if \( x(t) \cdot x^{\prime}(t) = 0 \) for all \( t \in I \), that is, if and only if \( x(t) \) is orthogonal to \( x^{\prime}(t) \) for all \( t \in I \).
\end{proof}

\begin{problem}{1.3.2}
	A circular disk of radius \( 1 \) in the plane \( xy \) rolls without slipping along the \( x \)-axis. The figure described by a point of the circumference of the disk is called a cycloid.
	\begin{enumerate}[label=(\alph*)]
		\item Obtain a parametrized curve \(\alpha: \mathbb{R} \rightarrow \mathbb{R}^2\) the trace of which is the cycloid, and determine its singular points.
		\item Compute the arc length of the cycloid corresponding to a complete rotation of the disc.
	\end{enumerate}
\end{problem}

\begin{problem}{1.3.4}
	Lete \(\alpha: (0, \pi) \rightarrow \mathbb{R}^2 \) be given by \(\alpha(t) = (\sin t, \cos t + \log \tan \frac{t}{2})\), where \( t \) is the angle that the \( y \)-axis makes with the vector \(\alpha^{\prime}(t)\). The trace of \(\alpha\) is called the tractix. Show that \begin{enumerate}[label=(\alph*)]
		\item \(\alpha\) is a differentiable parametrized curve, regular except at \( t = \frac{\pi}{2} \).
		\item The length of the segment of the tangent of the tractrix between the point of tangency and the \( y \)-axis is constantly equal to \( 1 \).
	\end{enumerate}
\end{problem}
\end{document}
